\documentclass[11pt]{article}
\usepackage[letterpaper,margin=1in]{geometry}
\usepackage[T1]{fontenc}
\usepackage{lmodern}
\usepackage{amsmath,amssymb,amsthm}
\usepackage{graphicx}
\usepackage{float}
\usepackage[colorlinks=true,linkcolor=blue,urlcolor=blue,citecolor=blue]{hyperref}
\usepackage{microtype}
\usepackage{enumitem}

\newtheorem{theorem}{Theorem}
\newtheorem{lemma}{Lemma}
\newcommand{\Rmet}{R}
\newcommand{\cD}{\mathcal D}
\newcommand{\Ball}{\operatorname{Ball}}

\title{Variable-Coefficient Order-Zero Operators on Compact Fractafolds\\
Are Calder\'on--Zygmund}
\author{Run \texttt{fa\_banach\_001} --- candidate full solution}
\date{August 11, 2026}

\begin{document}
\maketitle

\begin{abstract}
Ionescu, Rogers, and Strichartz conjectured that their
variable-coefficient order-zero pseudodifferential operators on compact
fractafolds are Calder\'on--Zygmund. We prove the conjecture in the standing
compact, boundaryless setting of their Section 9. Their eigenfunction
expansion reduces the kernel to an absolutely summable family of
constant-coefficient kernels. They summed only the size estimates; the
standard-kernel difference estimates are uniform over the same family and
sum identically. The obstruction that smooth functions need not be closed
under multiplication is therefore bypassed.
\end{abstract}

\section{The source conjecture}
Let $X$ be a compact fractafold without boundary in the setting of
Section 9 of Ionescu--Rogers--Strichartz \cite{IRS}. Write $\mu$ for its
measure, $\Rmet$ for the resistance metric, $d$ for its resistance-metric
dimension, and $-\Delta$ for its nonnegative Laplacian. A variable symbol
$p\in S^0$ satisfies, for all integers $j,r\ge0$,
\begin{equation}\label{eq:symbol}
 \sup_{x\in X,\lambda>0}
 \left| (\lambda\partial_\lambda)^r\Delta_x^j p(x,\lambda)\right|
 <\infty.
\end{equation}
The associated operator is
\[
 Tu(x)=p(x,-\Delta)u(x)
 =\sum_{\lambda\in\operatorname{spec}(-\Delta)}
 p(x,\lambda)P_\lambda u(x),\qquad u\in\cD.
\]
The source proves that $T$ is bounded on $L^2(\mu)$, has a continuous
off-diagonal kernel satisfying $|K(x,y)|\lesssim \Rmet(x,y)^{-d}$, and is
bounded on every $L^q$, $1<q<\infty$. It then conjectures that $T$ is a
Calder\'on--Zygmund operator in the sense of \cite{IR}.

In that cited definition, besides $L^2$ boundedness and the size estimate, a
kernel is standard if, for some $c>1$,
\begin{equation}\label{eq:standard}
 |K(x,y)-K(x,\bar y)|
 \lesssim
 \frac{\Rmet(y,\bar y)}{\Rmet(x,\bar y)}\,
 \Rmet(x,y)^{-d}
\end{equation}
whenever $\Rmet(x,y)\ge c\Rmet(y,\bar y)$. The displayed modulus is Dini.

\begin{figure}[H]
\centering
\includegraphics[width=0.88\textwidth]{figures/open_problem_crop.png}
\caption{The source's uniform kernel expansion and conjecture (arXiv:1108.2246,
PDF page 27).}
\end{figure}

\section{Idea of the proof}
The source expands $p$ in an eigenbasis in its $x$-variable. After multiplying
the $k$th coefficient by a high power of its eigenvalue, one obtains a bounded
family of constant-coefficient $S^0$ symbols. The associated kernels have
uniform Calder\'on--Zygmund estimates. Dividing back introduces scalar
coefficients whose sup norms are absolutely summable for a sufficiently high
power.

The source uses this decomposition only to sum kernel sizes. But the
Calder\'on--Zygmund difference is in the $y$-variable, while each expansion
coefficient depends only on $x$. Thus the difference estimates sum by the same
triangle inequality. No product of two functions is ever tested for membership
in $\operatorname{dom}\Delta$.

\section{Uniform constant-coefficient pieces}
Choose an orthonormal eigenbasis $(\phi_k)$,
\[
 -\Delta\phi_k=\lambda_k\phi_k,
 \qquad \lambda_k\ge0,
\]
and set
\[
 m_k(\lambda)=\int_X p(z,\lambda)\overline{\phi_k(z)}\,d\mu(z).
\]
We use $1+\lambda_k$ rather than $\lambda_k$ in order to include a possible
zero mode without a separate case.

\begin{lemma}\label{lem:uniform}
For every integer $N\ge0$, the symbols
\[
 b_{k,N}(\lambda)=(1+\lambda_k)^N m_k(\lambda)
\]
form a bounded family in $S^0$: every $S^0$ seminorm is bounded independently
of $k$.
\end{lemma}

\begin{proof}
Self-adjointness and the absence of boundary give
\begin{equation}\label{eq:ibp}
 b_{k,N}(\lambda)
 =\int_X (I-\Delta_z)^N p(z,\lambda)
     \overline{\phi_k(z)}\,d\mu(z).
\end{equation}
The operators $\lambda\partial_\lambda$ and $\Delta_z$ commute. Hence, for
every $r\ge0$, Cauchy--Schwarz and $\|\phi_k\|_2=1$ yield
\[
 |(\lambda\partial_\lambda)^r b_{k,N}(\lambda)|
 \le \mu(X)^{1/2}
 \max_{z\in X}|(I-\Delta_z)^N
       (\lambda\partial_\lambda)^r p(z,\lambda)|.
\]
Expanding $(I-\Delta_z)^N$ and using \eqref{eq:symbol} gives a bound depending
on $N,r,p$, but not on $k$ or $\lambda$.
\end{proof}

Let $B_{k,N}=b_{k,N}(-\Delta)$ and let $H_{k,N}(x,y)$ denote its kernel.
The constant-coefficient kernel theorem of \cite[Theorem 3.6]{IRS}, together
with the standard-kernel criterion of \cite{IR}, gives constants independent
of $k$ such that
\begin{align}
 |H_{k,N}(x,y)|&\le C_N\Rmet(x,y)^{-d},\label{eq:uniform-size}\\
 |H_{k,N}(x,y)-H_{k,N}(x,\bar y)|
 &\le C_N\frac{\Rmet(y,\bar y)}{\Rmet(x,\bar y)}
 \Rmet(x,y)^{-d}\label{eq:uniform-diff}
\end{align}
under the separation condition in \eqref{eq:standard}. Uniformity follows
from Lemma \ref{lem:uniform}: the proof of the constant-coefficient theorem
uses finitely many symbol seminorms for each displayed kernel estimate. In
particular it first proves, uniformly,
\[
 |\Delta_y H_{k,N}(x,y)|\le C_N\Rmet(x,y)^{-2d-1},
\]
from which \eqref{eq:uniform-diff} follows by the Ionescu--Rogers criterion.

\section{Summing the standard kernels}
\begin{theorem}\label{thm:main}
Under the standing compact, boundaryless hypotheses of Section 9 of
\cite{IRS}, every variable-coefficient order-zero operator $p(x,-\Delta)$ is a
Calder\'on--Zygmund operator in the precise sense cited by the source.
Consequently it is of weak type $(1,1)$ and bounded on $L^q(\mu)$ for every
$1<q<\infty$.
\end{theorem}

\begin{proof}
The eigenfunction estimate and compact spectral summability used in the proof
of \cite[Theorem 9.5]{IRS} allow us to choose $N$ so large that
\begin{equation}\label{eq:summability}
 A_N:=\sum_k
 \frac{\|\phi_k\|_\infty}{(1+\lambda_k)^N}<\infty.
\end{equation}
Indeed, the source proves $\|\phi_k\|_\infty\le C(1+\lambda_k)^\alpha$
for some $\alpha$ and then chooses an arbitrarily high spectral power with a
summable tail.

The eigenfunction expansion of $p(\cdot,\lambda)$ gives, first on finite
spectral sums,
\begin{align*}
 Tu(x)
 &=\sum_k \phi_k(x)m_k(-\Delta)u(x)\\
 &=\sum_k\frac{\phi_k(x)}{(1+\lambda_k)^N}B_{k,N}u(x).
\end{align*}
Thus the off-diagonal kernel found in the source is
\begin{equation}\label{eq:kernel-sum}
 K(x,y)=\sum_k
 \frac{\phi_k(x)}{(1+\lambda_k)^N}H_{k,N}(x,y).
\end{equation}
The series is locally uniformly convergent off the diagonal by
\eqref{eq:summability} and \eqref{eq:uniform-size}. Moreover,
\[
 |K(x,y)|
 \le C_N A_N\Rmet(x,y)^{-d}.
\]

For $\Rmet(x,y)\ge c\Rmet(y,\bar y)$, subtract \eqref{eq:kernel-sum} at
$y$ and $\bar y$, then use \eqref{eq:uniform-diff}:
\begin{align*}
 |K(x,y)-K(x,\bar y)|
 &\le \sum_k\frac{|\phi_k(x)|}{(1+\lambda_k)^N}
 |H_{k,N}(x,y)-H_{k,N}(x,\bar y)|\\
 &\le C_N A_N
 \frac{\Rmet(y,\bar y)}{\Rmet(x,\bar y)}
 \Rmet(x,y)^{-d}.
\end{align*}
This is \eqref{eq:standard} with a linear Dini modulus. The source already
proved the required $L^2$ boundedness and kernel representation. Hence $T$ is
Calder\'on--Zygmund. The endpoint and $L^q$ conclusions are the standard
consequences recorded in \cite{IR}; the strong $L^q$ part was also proved
directly in the source.
\end{proof}

\section{Verification, scope, and novelty}
There is a second route to the same result: the estimates above permit
termwise application of $\Delta_y$ in \eqref{eq:kernel-sum}; the resulting
series converges locally uniformly off the diagonal and is bounded by
$C A_N\Rmet(x,y)^{-2d-1}$. Applying the Ionescu--Rogers theorem directly then
recovers \eqref{eq:standard}.

The proof does not claim an extension to infinite blow-ups. In that setting
the compact eigenbasis and summability step \eqref{eq:summability} need a
replacement. It also uses the one-sided-in-$y$ standard-kernel definition
explicitly cited by the source; no assertion of smoothness in the $x$ variable
is made.

A bounded novelty search covered the run's registry, solution, attempt, and
proof-gap indexes; exact title and conjecture phrases; variable-coefficient
fractal pseudodifferential operators; and later citations. It found the source,
its published version, \cite{IR}, and an author research statement still
presenting the assertion as a conjecture, but no later resolution. Novelty is
therefore plausible but provisional.

\paragraph{Human-review recommendation.}
Review as a candidate full solution. The most sensitive point is the uniform
dependence of the source's constant-coefficient kernel estimates on a bounded
family of $S^0$ seminorms. This dependence is explicit in that proof, and the
source already uses it for the size estimate in its own variable-coefficient
expansion.

\begin{thebibliography}{9}
\bibitem{IRS}
M. Ionescu, L. G. Rogers, and R. S. Strichartz,
\emph{Pseudo-differential Operators on Fractals},
arXiv:1108.2246; Rev. Mat. Iberoam. 29 (2013), 1159--1190.

\bibitem{IR}
M. Ionescu and L. G. Rogers,
\emph{Complex Powers of the Laplacian on Affine Nested Fractals as
Calder\'on--Zygmund Operators},
arXiv:1002.2011.
\end{thebibliography}

\end{document}

